\section{Tipo societário escolhido e justificativas}
\label{sec:escolha}
Recomenda-se, para a atividade, a constituição do \projeto{} como \textbf{Sociedade Empresária Limitada (LTDA), com quatro sócios fundadores}. A escolha é uma proposta acadêmica coerente com a composição do grupo, e não declaração de empresa já constituída.

\subsection{Número de sócios e participação societária}
A distribuição abaixo é uma \textbf{hipótese igualitária de negociação} e não decorre de acordo já aprovado entre os fundadores. O contrato deverá registrar o capital efetivamente definido, sua forma e prazo de integralização e o número e valor das quotas (Código Civil, 2002).

\begin{center}
\begin{tabular}{@{}L{11cm} R{3cm}@{}}
\toprule
\cabfundo \cab{Sócio fundador proposto} & \cab{Participação proposta}\\
\midrule
Gabriel da Silva Cassino & 25\%\\
Giovanna Naves Ribeiro & 25\%\\
Júlia Rodrigues Vasconcellos Melo & 25\%\\
Kathleen Rodrigues Leite Barbosa & 25\%\\
\midrule
\totalfundo \textbf{Total} & \textbf{100\%}\\
\bottomrule
\end{tabular}
\end{center}

Não se fixa um capital social fictício. Seu valor deverá resultar do orçamento de abertura, da necessidade de caixa e dos aportes acordados. Trabalho futuro dos fundadores não pode ser usado como integralização de capital em uma limitada (art.~1.055, \S~2º, do Código Civil). Eventual aquisição gradual de participação por permanência e entregas, conhecida como \emph{vesting}, requer instrumento próprio, sem confusão com a integralização.

\subsection{Responsabilidade por dívidas e obrigações}
Nos termos do art.~1.052 do Código Civil, a responsabilidade de cada sócio é restrita ao valor de suas quotas, mas todos respondem solidariamente pela integralização do capital. Esse arranjo é mais adequado que o EI para uma plataforma que celebra contratos com clínicas, fornecedores e usuários.

A proteção exige separação entre contas pessoais e empresariais, escrituração e documentação dos repasses. Abuso da personalidade jurídica, desvio de finalidade e confusão patrimonial podem fundamentar desconsideração nos termos do art.~50. Garantias pessoais, atos ilícitos de administradores e outras hipóteses legais podem gerar responsabilização própria. Portanto, a LTDA não elimina riscos nem substitui gestão diligente (Código Civil, 2002).

\subsection{Facilidade de abertura e custos de manutenção}
A LTDA possibilita uma estrutura contratual menos complexa que a S.A., com registro na Junta Comercial e processo integrado pela Redesim. Exige despesas de registro, contabilidade, obrigações fiscais, eventual licenciamento e elaboração de instrumentos contratuais. Não se atribuem valores de taxas sem consulta atual à JUCEMG e aos profissionais responsáveis.

O planejamento deverá prever contabilidade, formalização e apoio jurídico/LGPD. Esses itens não comprovam que a abertura esteja paga nem que cubram todos os serviços. Antes da execução, deve-se separar custos únicos de abertura de custos recorrentes e obter propostas com escopo definido.

\subsection{Crescimento e atração de investidores}
A limitada permite admitir novos sócios e ajustar participações por alteração contratual. Acordo de sócios pode disciplinar preferência, critérios de avaliação, informação a investidores e diluição. A entrada de pessoa jurídica no capital pode impedir o Simples Nacional; assim, a rodada deve ser acompanhada de reavaliação tributária (LC nº 123/2006).

Se atendidos os requisitos do Marco Legal das Startups, podem ser avaliados instrumentos como opção de subscrição, mútuo conversível e investimento-anjo, com os limites e efeitos de cada modalidade. Não se deve tratar todo investidor como sócio automático. Uma transformação futura em S.A. poderá ser considerada quando houver necessidade concreta de capital e governança, sem adotá-la prematuramente (LC nº 182/2021).